\documentclass[11pt]{article}
\usepackage[T1]{fontenc}\usepackage[utf8]{inputenc}
\usepackage[margin=1.05in]{geometry}
\usepackage{amsmath,amssymb,amsthm,mathtools,lmodern,microtype}
\usepackage[authoryear,round]{natbib}\usepackage{xurl}
\usepackage[colorlinks=true,allcolors=blue]{hyperref}
\newcommand{\workbenchurl}{https://github.com/KokunoYumeto/schemes-and-moduli-workbench}
\newcommand{\recordurl}{https://github.com/KokunoYumeto/mathematics-workbenches}
\newcommand{\OO}{\mathcal O}\newcommand{\CC}{\mathcal C}
\newcommand{\A}{\mathbb A}\newcommand{\Gm}{\mathbb G_m}
\newcommand{\dR}{\mathrm{dR}}\newcommand{\et}{\mathrm{et}}
\DeclareMathOperator{\Spec}{Spec}\DeclareMathOperator{\Spwf}{Spwf}
\DeclareMathOperator{\Frac}{Frac}
\newtheorem{theorem}{Theorem}[section]
\newtheorem{lemma}[theorem]{Lemma}\newtheorem{proposition}[theorem]{Proposition}
\newtheorem{corollary}[theorem]{Corollary}
\theoremstyle{definition}\newtheorem{definition}[theorem]{Definition}
\theoremstyle{remark}\newtheorem{remark}[theorem]{Remark}
\title{Counterexamples to Grothendieck's weak-formal cohomology tests}
\author{}\date{}
\begin{document}\maketitle
\begin{abstract}
Grothendieck proposed a site of flat weak formal thickenings and asked whether it recovered Monsky--Washnitzer cohomology for smooth affine schemes and de Rham cohomology for smooth proper liftable schemes. We calculate its first structural-sheaf cohomology on the affine line and the multiplicative group: both groups vanish. The multiplicative group therefore contradicts the affine comparison even after tensoring cohomology with the fraction field. We also prove vanishing for every smooth proper lift whose special fibre has no first \'etale cohomology with $\mathbf F_p$ coefficients. This uses Ogus's finite-level infinitesimal comparison and Meredith's weak formal cohomology theorem. A supersingular elliptic curve gives a proper counterexample, with de Rham cohomology of rank two. The distinction between tensoring cohomology and rationalizing its coefficient sheaf is essential. The mathematics and text were produced by GPT-6 Astra Pro in ChatGPT Pro mode. A separate statement-only reconstruction of all main results was frozen and compared with these proofs; no human has checked them.
\end{abstract}

\section{Grothendieck's two tests}

The choice of thickenings and coefficients matters in a site intended to recover $p$-adic cohomology. Grothendieck's manuscript, cote~7, defines a weak-formal site on archive sheet~17 and sets out its decisive tests on sheet~18. The latter sheet asks for the Monsky--Washnitzer definition on a smooth affine scheme and, for a smooth proper lift, says
\begin{quote}
``on doit retrouver (du moins mod torsion) la cohomologie de De Rham''
\end{quote}
\citep[archive sheet~18; scan image~19]{Grothendieck}.
Thus the proper test explicitly allows discarding torsion. We show that the site defined on the preceding sheet fails this test even after that allowance. We also show a rational failure of its affine test.

Let $k$ be a perfect field of characteristic $p>0$, let $W=W(k)$ be its ring of Witt vectors, and let $K=\Frac W$. An integral dagger algebra is a quotient of
\[
 W\langle x_1,\ldots,x_r\rangle^\dagger
 =\left\{\sum_{\nu}a_\nu x^\nu:
       v_p(a_\nu)\geq c|\nu|-C
       \text{ for some }c>0,C\right\}.
\]
It is required to be $p$-adically separated. Such algebras are also called weakly complete finitely generated $W$-algebras. Their associated weak formal schemes have underlying space $\Spec(B/pB)$ and dagger localizations as their local rings. We use integral dagger algebras throughout.

\begin{definition}\label{def:site}
For a $k$-scheme $X$ locally of finite type, the site $\CC_X$ has objects
\[
 (U,\Spwf B),\qquad
 U\subset X\text{ affine open},\quad
 B/pB\twoheadrightarrow\OO_X(U),
\]
where $B$ is a flat integral dagger $W$-algebra and the displayed kernel is nilpotent. Morphisms are the compatible morphisms of weak formal schemes and open subschemes of $X$. Coverings are Zariski open coverings of the weak formal scheme. The structural sheaf $\OO$ has value $B$ on this object. There are no divided powers among the data.
\end{definition}

This is the flat weak-formal site on archive sheet~17, scan image~18, of the cited manuscript. It is not a definition of the crystalline site. Its objects are also more restrictive than the enlargements in a modern overconvergent site.

\begin{theorem}\label{thm:affine}
For every perfect field $k$ of positive characteristic,
\[
 H^1(\CC_{\A^1_k},\OO)=0,
 \qquad H^1(\CC_{\mathbb G_{m,k}},\OO)=0.
\]
In contrast, integral Monsky--Washnitzer cohomology satisfies
\[
 [x^{p-1}dx]\ne0\quad\text{in }H^1_{\rm MW}(\A^1_k/W),
 \qquad p[x^{p-1}dx]=0,
\]
and
\[
 H^1_{\rm MW}(\mathbb G_{m,k}/W)\otimes_WK=K[dx/x].
\]
The affine comparison therefore fails integrally on $\A^1_k$ and after tensoring cohomology with $K$ on $\mathbb G_{m,k}$.
\end{theorem}

\begin{theorem}\label{thm:proper}
Let $k$ be algebraically closed of characteristic $p>0$, and let $\mathcal X$ be a nonempty smooth proper $W(k)$-scheme with geometrically connected fibres. Write $X=\mathcal X_k$. If
\[
 H^1_{\et}(X,\mathbf F_p)=0,
\]
then $H^1(\CC_X,\OO)=0$.
\end{theorem}

\begin{corollary}\label{cor:elliptic}
Let $k=\overline{\mathbf F}_3$ and let $\mathcal E/W(k)$ be the smooth projective elliptic curve
\[
 y^2z=x^3-xz^2.
\]
For its special fibre $E$,
\[
 H^1(\CC_E,\OO)=0,
 \qquad H^1_{\dR}(\mathcal E/W(k))\simeq W(k)^2.
\]
Thus Grothendieck's proper test fails even modulo torsion.
\end{corollary}

These statements concern $H^1(\CC_X,\OO)$ and its tensor product with $K$. They make no assertion that
\[
 H^1(\CC_X,\OO)\otimes_WK
       =H^1(\CC_X,\OO[1/p]).
\]
Such an interchange is not part of the argument. The failure established here does not apply to a cohomology theory with a different site or coefficient sheaf.

\subsection*{Earlier results and the mechanism of the proofs}

\citet[Corollary~4.7, p.~314]{Ogus} proves that, for a proper scheme over an algebraically closed field which embeds in a smooth $W_n(k)$-scheme, its ordinary infinitesimal cohomology with structural coefficients is its \'etale cohomology with constant $W_n(k)$ coefficients. This is a finite-level theorem, without divided powers. It supplies the vanishing used in Theorem~\ref{thm:proper}; we do not claim that theorem as a new infinitesimal comparison. Our argument explains why its degree-one vanishing detects torsors on the specified flat weak-formal site.

\citet[\S2, Theorem~14, p.~10]{Meredith} proves the affine weak-formal acyclicity needed to trivialize torsors on a smooth affine lift. His proper comparison, \S5, Theorems~A and~B, p.~25, identifies coherent cohomology with that of the weak completion. We use the cohomology comparison for proper schemes; we do not require an algebraization equivalence for all coherent sheaves on an arbitrary proper weak formal scheme. \citet[Proposition~2.31, author version p.~19]{DLZ} gives the Hensel property for weakly complete algebras that permits the local smooth lifts used below.

The integral Monsky--Washnitzer group in Theorem~\ref{thm:affine} does not depend on an arbitrary choice of lift. \citet[Proposition~2 and Lemma~3, author version pp.~4--5]{ErtlSprang} prove canonical independence and functoriality without killing torsion. Their Proposition~1, pp.~3--4, proves the stronger input: maps from a smooth weak lift to any weakly complete $W$-algebra which agree modulo $p$ induce chain-homotopic maps on integral continuous de Rham complexes. The homotopy integrates between $0$ and $p$, with coefficients $p^{i+1}/(i+1)$, which are integral and converge rapidly. These are known results; the torsion class displayed here is a direct calculation in that canonical theory.

There is an earlier failure theorem for a related dagger site. \citet[Theorem~8.6, pp.~1987--1989, and Corollary~8.11, p.~1989]{ArabiaMebkhout} construct a nonzero first cohomology class on affine space with rational structural coefficients. Their Definition~3.1, p.~1934, requires each flat dagger lift to reduce exactly to the given open subscheme; its kernel is $p\OO$. Definition~\ref{def:site} also permits nilpotent thickening ideals in $B/pB$, including the diagonal tests used below. Moreover our vanishing is for the integral sheaf before tensoring its cohomology with $K$. The two conventions differ in both objects and coefficients. Their theorem already establishes that a dagger infinitesimal structural sheaf need not have the expected de Rham cohomology. The results here calculate the particular enlarged nilpotent-test site in the manuscript and establish the proper vanishing criterion.

\citet[Definitions~3.1 and~3.7, Example~3.8(ii)]{Zhang} constructs a Monsky--Washnitzer topos using enlargements with arbitrary special-fibre maps to $X$. Its analytic structural sheaf has value $B[1/p]$. His Example~4.4, p.~17, includes universal enlargements
\[
 W\langle x,s\rangle^\dagger/(x^n-ps).
\]
Over a point their reduction-to-$k$ kernel contains $s$, which is not nilpotent. They are therefore absent from Definition~\ref{def:site}. His Corollary~5.19, p.~27, gives the analytic rigid-cohomology comparison for that larger site and its rational coefficient sheaf. His Example~3.8(ii), p.~14, explicitly distinguishes rationalizing a sheaf from tensoring its cohomology. These differences explain why the counterexamples here do not contradict that comparison.

\citet[\S3, pp.~834--835]{Illusie} reviews the crystalline comparison and its divided-power Poincar\'e lemma, and explains how the same comparison is expressed through recent de Rham stacks. That account concerns divided-power thickenings and their de Rham comparison. Our counterexamples concern ordinary nilpotent thickenings with the flat weak-formal coefficient convention in Definition~\ref{def:site}. No assertion about the crystalline comparison is altered here.

The main difficulty in the affine calculation is proving that calculations on nilpotent diagonal neighbourhoods detect a cocycle on every permitted weak formal thickening. We do this by a finite weak Nakayama argument and a flat fibre product. On the multiplicative group, integrality then excludes the logarithm, whose derivative is $dx/x$. In the proper case, Ogus's theorem removes every finite-level stratified extension. Proper coherent cohomology and $p$-adic separatedness remove the remaining extension and its cocycle. The same flat fibre product finishes the passage to arbitrary objects.

\section{Finite weak algebras and detection of cocycles}

This section gives the algebraic step which ensures that calculations on infinitesimal diagonal neighbourhoods control the entire site.

\begin{lemma}\label{lem:finite}
Let $A\to B$ be a homomorphism of integral dagger $W$-algebras. If $B/pB$ is finite over $A/pA$, then $B$ is a finite $A$-module.
\end{lemma}
\begin{proof}
Choose finitely many weak algebra generators $y_1,\ldots,y_r$ of $B$ over $A$. Integrality modulo $p$ gives a monic polynomial $P_j(T)\in A[T]$ of degree $d_j$ with $P_j(y_j)\in pB$. Choose a dagger series $G_j$ representing $P_j(y_j)/p$ in the presentation by the $y_j$. Then $B$ is a quotient of
\[
 A\langle y_1,\ldots,y_r\rangle^\dagger/
        (P_j(y_j)-pG_j:1\leq j\leq r).
\]
We prove that the monomials $y^\nu$ with $0\leq\nu_j<d_j$ generate this quotient as an $A$-module.

Present $A$ as a quotient of an integral dagger algebra in finitely many base variables and lift all coefficients to that presentation. Choose radii greater than one at which the finitely many coefficients of the $P_j$ and the series $G_j$ converge. The radii can be chosen arbitrarily close to one. Choose each $y_j$ radius $\rho_j$ large enough relative to the base radii that every lower term of $P_j$ has norm at most $\rho_j^{d_j}$, and choose all radii still within the convergence range. Successive monic division is then bounded in the corresponding Gauss norms. Its remainder operator $R$ has norm at most one, and its quotient operators $Q_j$ have norm at most $\rho_j^{-d_j}$. These inequalities follow by replacing $y_j^{d_j}$ by the negative lower terms: each replacement weakly decreases the weighted norm and lowers that exponent.

As the radii approach one, the norms of the integral $G_j$ approach values at most one. We may consequently arrange that
\[
 T(f)=p\sum_jG_jQ_j(f)
 \quad\text{has operator norm less than one}.
\]
The division identity becomes
\[
 f=R(f)+\sum_j(P_j-pG_j)Q_j(f)+T(f).
\]
Iterating it yields a convergent remainder
\[
 \sum_{a\geq0}R(T^af)
\]
with bounded $y$-degrees. Its coefficients converge in a fixed base Banach algebra at radii greater than one. The quotient series converge there too, so the resulting identity is an identity of dagger series, not only of completed series. All coefficients are integral: monic division by integral polynomials preserves integrality, and $T^a$ introduces $p^a$. Thus the remainder coefficients belong to the integral algebra $A$. The stated finitely many monomials generate, proving the lemma.
\end{proof}

We will use the following precise detection statement. A natural additive cocycle assigns $u_B(a_0,a_1)\in B$ to pairs of compatible maps from a smooth affine lift, functorially in the object $B$, and satisfies
\[
 u_B(a_0,a_2)=u_B(a_0,a_1)+u_B(a_1,a_2).
\]
Such a cocycle describes a structural-sheaf torsor once it has been trivialized on the lift.

\begin{lemma}\label{lem:detect}
Let $A$ be a smooth integral dagger lift of an affine open $U\subset X$. Suppose a natural additive cocycle vanishes on all finite diagonal neighbourhoods of this lift and their open restrictions. Then it vanishes on every object over $U$ and every pair of compatible maps $A\rightrightarrows B$.
\end{lemma}
\begin{proof}
Write $I$ for the diagonal ideal and
\[
 P_N=(A\mathbin{\widehat\otimes^\dagger_W}A)/I^N.
\]
Here the dagger tensor notation denotes weak completion of the algebraic tensor product. The quotient is a finite $A$-module. Since the lift is smooth, its diagonal is regular and the successive quotients of its ideal filtration are symmetric powers of the locally free module $\Omega^1_{A/W}$. Hence $P_N$ is flat over $W$. It is an allowed object over $U$.

For a pair $a_0,a_1:A\rightrightarrows B$, the image of $I$ in $B/pB$ is contained in the nilpotent thickening ideal. For each $m\geq1$ some $N$ therefore satisfies $I^N B\subset p^mB$. The pair induces $P_N\to B/p^mB$. Form the ordinary ring fibre product
\[
 C=B\times_{B/p^mB}P_N.
\]
Both factors are finite over $A$, using the first projection for the $A$-action. For $B$, choose finitely many algebra generators of $B/pB$ over $A/pA$. Each differs from an element lifted from $\OO(U)$ by a nilpotent element, since the composite $A/pA\to B/pB\to\OO(U)$ is the identity. These generators are integral over $A/pA$. Lemma~\ref{lem:finite} now applies. The finiteness of $P_N$ was just observed. Integral dagger algebras are Noetherian, so the submodule $C\subset B\oplus P_N$ is also finite over $A$.

The ring $C$ is weakly complete. Indeed a dagger series evaluated componentwise stays in the fibre product: modulo $p^m$ it has only finitely many nonzero terms, and its two evaluations agree there. A finite set of $A$-module generators, together with weak generators for $A$, makes $C$ a dagger algebra of finite type over $W$. It is $p$-adically separated and $W$-flat as a submodule of the separated flat module $B\oplus P_N$.

The map $C\to P_N$ is onto, with kernel $(p^mB,0)$. Modulo $pC$ this ideal is square-zero, since
\[
 (p^mb,0)(p^mb',0)
       =p(p^{2m-1}bb',0),\qquad 2m-1\geq m.
\]
The kernel of $P_N/pP_N\to\OO(U)$ is nilpotent. It follows that $C/pC\to\OO(U)$ is onto with nilpotent kernel. Thus $C$ is an allowed object of the original site; no divided powers or additional enlargements have been introduced.

The pair of maps into $B$ and the universal pair into $P_N$ combine into a pair $A\rightrightarrows C$. The value of the cocycle on $C$ maps to zero in $P_N$ by hypothesis. Its image in $B$ consequently belongs to $p^mB$. This is true for every $m$, so separatedness gives $u_B(a_0,a_1)=0$.
\end{proof}

For a general smooth lift this argument is local on $U$. Maps into every test object exist locally: lift smooth coordinates through its nilpotent special-fibre ideal and lift the \'etale part using the Hensel property of the dagger algebra \citep[Proposition~2.31, p.~19]{DLZ}. Equivalently one uses a standard smooth presentation and the invertibility of its Jacobian minor. The Hensel property gives a solution in the weak algebra itself. A solution only in its completion would not suffice.

\section{The affine line and the multiplicative group}

We calculate the torsors by their Taylor cocycles and retain the integral condition at every coefficient.

Put
\[
 A_+=W\langle x\rangle^\dagger,
 \qquad A_\times=W\langle x,x^{-1}\rangle^\dagger.
\]
Their represented sheaves cover the terminal sheaf on the corresponding sites. For the affine line, a lift of its coordinate gives a map to every object. For the multiplicative group, that lift is a unit: it is a unit modulo the nilpotent ideal in $B/pB$, and $p$ belongs to the Jacobson radical of a dagger algebra. The latter assertion follows, for example, by the convergent geometric series for $1-pb$.

An additive $\OO$-torsor becomes trivial on either of these objects. Its local trivializations there use Zariski weak formal open coverings, and affine structural-sheaf acyclicity is Meredith's Theorem~14 \citep[p.~10]{Meredith}. Choosing such a trivialization gives a natural cocycle $u_B(a_0,a_1)$. A change of trivialization subtracts
\[
 a_1(f)-a_0(f),\qquad f\in A_+\text{ or }A_\times.
\]

First work with $A_\times$. Evaluate the cocycle on the allowed objects
\[
 P_N=A_\times[t]/(t^N),
 \qquad a_0(x)=x,\quad a_1(x)=x+t.
\]
Compatibility in $N$ gives a series $u(x,t)\in A_\times[[t]]$ with zero constant term. Evaluation on the two-variable diagonal neighbourhoods gives
\begin{equation}\label{eq:cocycle}
 u(x,t+s)=u(x,t)+u(x+t,s).
\end{equation}
Define $a(x)=\partial_tu(x,0)\in A_\times$. Differentiating \eqref{eq:cocycle} in $s$ at zero shows that
\[
 \partial_tu(x,t)=a(x+t).
\]
Write
\[
 a(x)=\sum_{j\in\mathbb Z}a_jx^j,\qquad
 c=a_{-1},\qquad b_r=a_{r-1}/r\quad(r\ne0).
\]
Formal integration over $K$ gives the unique solution with zero constant term:
\begin{equation}\label{eq:log}
 u(x,t)=c\log(1+t/x)
       +\sum_{r\ne0}b_r\big((x+t)^r-x^r\big).
\end{equation}
Each coefficient is a well-defined overconvergent Laurent series. Taking finitely many derivatives costs at most logarithmic valuation in the exponent, which does not destroy linear overconvergence.

All coefficients in \eqref{eq:log} are integral because $u\in A_\times[[t]]$. The coefficient of $t^nx^{-n}$ is $(-1)^{n-1}c/n$: no term indexed by $r\ne0$ contributes to this Laurent exponent. Taking $n=p^a$ for arbitrarily large $a$ forces $c=0$.

Fix $r\ne0$. The coefficient of $t^nx^{r-n}$ is $b_r\binom rn$. Set $n=p^{v_p(r)}$. Then $\binom rn$ is a $p$-adic unit. This also holds when $r$ is negative: in $\mathbf F_p[[z]]$, write $r=p^av$ with $p\nmid v$ and use
\[
 (1+z)^r=(1+z^{p^a})^v.
\]
Its first nonconstant coefficient is $v$ at degree $p^a$. Thus $b_r\in W$ for every nonzero $r$.

There are constants $\epsilon>0$ and $C$ such that $v_p(a_j)\geq\epsilon|j|-C$. Since $v_p(r)\leq\log_p|r|$, division by $r$ leaves a linear lower bound of positive slope outside a finite set. It follows that
\[
 f(x)=\sum_{r\ne0}b_rx^r\in A_\times.
\]
Subtracting its coboundary kills the cocycle on every $P_N$.

To apply Lemma~\ref{lem:detect} over smaller opens, take a principal open with dagger localization $A_U$. Both maps from $A_\times$ extend to $A_U$, since the chosen denominator is a unit modulo the nilpotent thickening ideal and hence a unit. The vanishing on its finite diagonal neighbourhoods follows by localization. The lemma shows that the cocycle vanishes on every object. The torsor is trivial and $H^1(\CC_{\mathbb G_{m,k}},\OO)=0$.

For $A_+$ the proof is the same with $a(x)=\sum_{j\geq0}a_jx^j$. There is no logarithmic term and only positive indices $r=j+1$. The same unit binomial coefficient forces $b_r\in W$, the same decay estimate puts the primitive in $A_+$, and Lemma~\ref{lem:detect} kills the remaining cocycle. Hence $H^1(\CC_{\A^1_k},\OO)=0$ as well.

For comparison, the integral Monsky--Washnitzer complex of the affine line contains the nonzero class $x^{p-1}dx$. Its multiple by $p$ is $d(x^p)$, whereas an integral primitive would have coefficient $1/p$ at $x^p$, which is impossible. On the multiplicative group over $K$, termwise integration removes every Laurent differential except its residue times $dx/x$. The same decay estimate ensures that the primitive remains overconvergent. A derivative has zero residue, so $dx/x$ is not exact. This proves all assertions of Theorem~\ref{thm:affine}.

\begin{remark}\label{rem:torsion}
The absence of $p$-torsion on the affine-line site also has a short direct explanation. Flatness gives
\[
 0\longrightarrow\OO\xrightarrow{p}\OO
     \longrightarrow\OO/p\longrightarrow0.
\]
A global section of $\OO/p$ is determined by a polynomial $f(x)\in k[x]$ on the smooth affine lift. The two maps into $(W[x,t]/t^N)^\dagger$ force $f(x+t)=f(x)\bmod t^N$ for every $N$, so $f$ is constant. Thus $H^0(\CC_{\A^1_k},\OO/p)=k$, and its elements lift through $W$. The long exact sequence gives $H^1(\CC_{\A^1_k},\OO)[p]=0$. This simpler observation already separates the site from integral de Rham cohomology; the full torsor calculation proves the stronger vanishing above.
\end{remark}

\section{Proper schemes with vanishing first \'etale cohomology}

We now prove Theorem~\ref{thm:proper}, using finite-level infinitesimal cohomology to control the stratification of a weak-formal torsor.

Write $W_n=W/p^nW$ and $\mathcal X_n=\mathcal X\times_WW_n$. From the hypothesis and the coefficient sequences
\[
 0\longrightarrow\mathbb Z/p^{n-1}\mathbb Z
 \longrightarrow\mathbb Z/p^n\mathbb Z
 \longrightarrow\mathbb Z/p\mathbb Z\longrightarrow0
\]
we obtain by induction
\[
 H^1_{\et}(X,\mathbb Z/p^n\mathbb Z)=0\quad(n\geq1).
\]
The embedding $X\hookrightarrow\mathcal X_n$ is into a smooth $W_n$-scheme. Ogus's Corollary~4.7, p.~314, therefore gives
\begin{equation}\label{eq:infzero}
 H^1((X/W_n)_{\rm inf},\OO)=0\quad(n\geq1).
\end{equation}
This use of his comparison satisfies its properness, algebraically closed residue field, and smooth-embedding hypotheses at every finite level.

Let $\mathcal T$ be an additive $\OO$-torsor on $\CC_X$. Restricting to the affine opens of the weak completion $\mathcal X^\dagger$ and gluing gives an extension
\begin{equation}\label{eq:extension}
 0\longrightarrow\OO_{\mathcal X^\dagger}
 \longrightarrow V\longrightarrow\OO_{\mathcal X^\dagger}
 \longrightarrow0.
\end{equation}
It is locally free of rank two: locally the quotient is free and the sequence splits. This construction identifies additive torsors with extensions of the indicated form, with fixed subobject and quotient.

Let $P_N$ be the scheme cut out by the $N$th power of the diagonal ideal in $\mathcal X\times_W\mathcal X$. Its projection to $\mathcal X$ is finite. Its successive diagonal-ideal quotients are symmetric powers of $\Omega^1_{\mathcal X/W}$; it is consequently $W$-flat. The two pullbacks of \eqref{eq:extension} to $P_N^\dagger$ are canonically identified by the torsor. These identifications are the identity on the diagonal, preserve the subobject and quotient, and satisfy the ordinary triple cocycle. Thus $V$ has an ordinary stratification on every finite diagonal neighbourhood.

Reduce this stratified extension modulo $p^n$. It defines a coherent crystal extension on $(X/W_n)_{\rm inf}$. Here is the descent construction explicitly. Every nilpotent test thickening of an open of $X$ admits locally a map to $\mathcal X_n$, by smoothness. Pull back the vector bundle along that map. Two such maps have differences in the nilpotent ideal, so their pair factors through a finite diagonal neighbourhood; its stratification identifies the two pullbacks. The triple identity makes these identifications independent of the choice and permits gluing. This also constructs the extension maps and preserves their exactness, since the quotient is locally free. A horizontal splitting of the crystal is exactly a splitting of the extension compatible with all finite diagonal identifications.

By \eqref{eq:infzero}, this crystal extension splits. Indeed extensions of the free rank-one structural module by itself in a ringed topos are classified by $H^1$ of its structural sheaf. We conclude that \eqref{eq:extension} has a horizontal splitting modulo every $p^n$.

Meredith's proper cohomology comparison identifies the underlying extension class with an element of the finite $W$-module
\[
 H^1(\mathcal X^\dagger,\OO)
       \simeq H^1(\mathcal X,\OO_{\mathcal X}).
\]
Its reduction modulo every $p^n$ is zero. The exact sequence for multiplication by $p^n$ on $\OO_{\mathcal X}$ shows that the class belongs to $p^nH^1(\mathcal X,\OO)$ for every $n$. A finite module over the complete discrete valuation ring $W$ is $p$-adically separated. Hence the extension class is zero. Fix a global splitting of \eqref{eq:extension} on $\mathcal X^\dagger$.

In this splitting, its stratification is an additive cocycle
\[
 u_N\in H^0(P_N^\dagger,\OO)
\]
vanishing on the diagonal. We show that each $u_N$ is zero. First,
\begin{equation}\label{eq:constants}
 H^0(\mathcal X_n,\OO)=W_n\quad(n\geq1).
\end{equation}
For $n=1$ this is proper geometric connectedness and smoothness over the algebraically closed field. Inductively the sequence
\[
 0\longrightarrow\OO_X\xrightarrow{p^{n-1}}
 \OO_{\mathcal X_n}\longrightarrow\OO_{\mathcal X_{n-1}}
 \longrightarrow0
\]
shows that a global section reducing to a constant differs from a lift of that constant by an element of $H^0(X,\OO_X)=k$. Constants lift through $W_n$, so this proves surjectivity in \eqref{eq:constants}. Injectivity follows from flatness and the nonempty special fibre.

The chosen splitting modulo $p^n$ and a horizontal splitting differ by a global function on $\mathcal X_n$. By \eqref{eq:constants} this function is constant. Its diagonal coboundary is zero. Therefore $u_N$ reduces to zero modulo every $p^n$.

The scheme $P_N$ is proper over $W$ and $W$-flat. Meredith's comparison identifies $H^0(P_N^\dagger,\OO)$ with the finite $W$-module $H^0(P_N,\OO)$. The multiplication exact sequence implies that a section vanishing modulo $p^n$ lies in $p^nH^0(P_N,\OO)$. Separatedness now gives $u_N=0$.

Finally take any old-site object and work over an affine open of $X$ with a smooth affine lift inside $\mathcal X$. Local maps to that lift exist by the Hensel argument following Lemma~\ref{lem:detect}. In the global splitting just chosen, all finite diagonal cocycles vanish. That lemma makes the cocycle vanish on every test object. The splitting consequently defines a section of $\mathcal T$ on the entire site. The torsor is trivial. This proves Theorem~\ref{thm:proper}.

\section{The elliptic counterexample}

The example exhibits both sides of the failed proper comparison by direct calculations.

The Weierstrass discriminant of $y^2=x^3-x$ is $64$, a $3$-adic unit, so its projective model is a smooth proper elliptic curve over $W(\overline{\mathbf F}_3)$. On the special fibre, a nonzero invariant differential is $dx/y$ up to a nonzero scalar. In characteristic three the equation gives
\[
 \frac{dx}{y}
 =\left(\frac{x}{y}\right)^3dx
   -\left(\frac1y\right)^3x\,dx.
\]
The Cartier operator kills both $dx$ and $x\,dx$, since it selects the exponents congruent to $2$ modulo $3$. Its semilinearity therefore kills the displayed invariant differential. By Serre duality, Frobenius on $H^1(E,\OO_E)$ is zero.

The Artin--Schreier sequence is exact on the \'etale site:
\[
 0\longrightarrow\mathbf F_3\longrightarrow\OO_E
    \xrightarrow{F-1}\OO_E\longrightarrow0.
\]
On $H^0(E,\OO_E)=k$, the map $a\mapsto a^3-a$ is onto because $k$ is algebraically closed. On $H^1(E,\OO_E)$ the map $F-1$ is minus the identity. Its kernel is zero, and hence $H^1_{\et}(E,\mathbf F_3)=0$. Theorem~\ref{thm:proper} gives $H^1(\CC_E,\OO)=0$.

For the lift, the degree-one Hodge sequence is
\[
 0\longrightarrow H^0(\mathcal E,\Omega^1_{\mathcal E/W})
 \longrightarrow H^1_{\dR}(\mathcal E/W)
 \longrightarrow H^1(\mathcal E,\OO_{\mathcal E})
 \longrightarrow0.
\]
Both outer modules are free of rank one. To verify that there is no intervening differential, $d$ vanishes on global constants. For the remaining differential $H^1(\OO)\to H^1(\Omega^1)$, flat base change identifies its generic fibre with the corresponding map on a smooth proper curve over $K$. That field-level map is zero by \citet[Lemma~50.20.3, Tag~0FW6]{Stacks}. Its integral target is torsion-free, so the integral map is zero as well. Thus the displayed sequence is exact integrally. It splits as a sequence of $W$-modules because its right term is free. Its middle term is $W^2$, which proves Corollary~\ref{cor:elliptic}.

\section{Scope and remaining questions}

The affine and proper examples settle the two tests for Definition~\ref{def:site} with integral structural coefficients, including the interpretation obtained by tensoring their cohomology with $K$. Theorem~\ref{thm:proper} is a vanishing criterion, not a general identification of this weak-formal cohomology with \'etale cohomology. Its proof uses a smooth proper lift, algebraic closedness of the residue field, and the vanishing of first \'etale cohomology. Removing any of these hypotheses requires a separate argument.

Two precise questions remain beyond these results. First, determine $H^1(\CC_X,\OO)$ for a general smooth proper lift with nonzero $H^1_{\et}(X,\mathbf F_p)$, including which finite-level infinitesimal classes lift compatibly to the flat weak-formal site. The proof above only detects the case where all those classes vanish. Second, compute $H^1(\CC_X,\OO[1/p])$ on the same restricted site and compare it with $H^1(\CC_X,\OO)\otimes_WK$. The integral Taylor coefficient argument cannot be used unchanged after rationalizing the coefficient sheaf. Neither Arabia--Mebkhout's smaller lifting site nor Zhang's larger enlargement site is the restricted nilpotent-test site here, so their results do not answer this latter question by themselves.

\section*{Note on method and reuse}

The mathematics and text were produced by GPT-6 Astra Pro in ChatGPT Pro mode. A separate instance of the same model received the definitions and main statements, reconstructed all three main results, and froze its proofs before comparison. That reconstruction uses finite interpolation to detect every allowed test and Frobenius-divided extensions on a proper lift. Its full proofs and elliptic calculation were compared with this article; prior exposure and unsuccessful approaches are recorded. The auxiliary finite weak Nakayama lemma has a targeted check only. Reasoning-effort metadata is unavailable. No human has checked the mathematics or text. The text is dedicated to the public domain under CC0~1.0. Sources and records are at \url{\workbenchurl} and \url{\recordurl}.

\small
\begin{thebibliography}{9}
\bibitem[Arabia and Mebkhout(2010)]{ArabiaMebkhout}
Arabia, A., \& Mebkhout, Z. (2010). Sur le Topos infinit\'esimal $p$-adique d'un sch\'ema lisse I. \emph{Annales de l'Institut Fourier, 60}(6), 1905--2094. \url{https://doi.org/10.5802/aif.2576}.
\bibitem[Davis et~al.(2012)]{DLZ}
Davis, C., Langer, A., \& Zink, T. (2012). Overconvergent Witt vectors. \emph{Journal f\"ur die reine und angewandte Mathematik, 668}, 1--34. \url{https://doi.org/10.1515/CRELLE.2011.141}. Author version: \url{https://arxiv.org/abs/1008.0305}.
\bibitem[Ertl and Sprang(2018)]{ErtlSprang}
Ertl, V., \& Sprang, J. (2018). Integral comparison of Monsky--Washnitzer and overconvergent de Rham--Witt cohomology. \emph{Proceedings of the American Mathematical Society, Series B, 5}, 64--72. \url{https://doi.org/10.1090/bproc/38}. Author version: \url{https://arxiv.org/abs/1709.07315v2}.
\bibitem[Grothendieck(n.d.)]{Grothendieck}
Grothendieck, A. (n.d.). \emph{Manuscript, cote~7} [Unpublished manuscript]. Archives Grothendieck, Universit\'e de Montpellier. Archive sheets~17--18, scan images~18--19. \url{https://grothendieck.umontpellier.fr/7.pdf}.
\bibitem[Illusie(2025)]{Illusie}
Illusie, L. (2025). Grothendieck and differential calculus. \emph{Comptes Rendus Math\'ematique, 363}, 829--841. \url{https://doi.org/10.5802/crmath.764}.
\bibitem[Meredith(1972)]{Meredith}
Meredith, D. (1972). Weak formal schemes. \emph{Nagoya Mathematical Journal, 45}, 1--38. \url{https://doi.org/10.1017/S0027763000014641}.
\bibitem[Ogus(1975)]{Ogus}
Ogus, A. (1975). Cohomology of the infinitesimal site. \emph{Annales scientifiques de l'\'Ecole normale sup\'erieure, 8}(3), 295--318. \url{https://doi.org/10.24033/asens.1289}.
\bibitem[The Stacks Project Authors(2026)]{Stacks}
The Stacks Project Authors. (2026). \emph{The Stacks Project}. Lemma~50.20.3, Tag~0FW6. \url{https://stacks.math.columbia.edu/tag/0FW6}.
\bibitem[Zhang(2019)]{Zhang}
Zhang, D. (2019). \emph{The Monsky--Washnitzer topos} [Preprint]. arXiv. \url{https://arxiv.org/abs/1904.09272}.
\end{thebibliography}
\end{document}
